%======================================================================
\chapter{Separation by logarithms; Janiszewski and Borsuk}\label{ch:sep}
%======================================================================

This chapter relates the separation properties of a plane set $K$ to logarithms on $K$. The two
basic facts are: if $y$ lies in the component of $x$ in $\R^2\setminus K$, then
$(z-x)/(z-y)$ has a logarithm on $K$ (Theorem~\ref{thm:log-same-comp}); if the component of $x$ is
bounded and does not contain $y$, it has none (Theorem~\ref{thm:nolog-bounded}). From them follow
Janiszewski's theorem (Theorem~\ref{thm:janiszewski}), Borsuk's theorem in the direction ``a compact
set with connected complement carries logarithms of all nonvanishing functions''
(Theorem~\ref{thm:borsuk}), and the topological invariance of plane separation
(Theorem~\ref{thm:sep-invariant}). None of these uses the Jordan curve theorem.

\section{Complementary components}

\begin{definition}\label{def:separates}
Let $A\subseteq\R^2$. For $x\notin A$ write $\comp Ax$ for the component of $x$ in $\R^2\setminus A$.
\begin{itemize}
\item $A$ \emph{separates the plane} if $\R^2\setminus A$ is not connected.
\item $A$ \emph{separates $x$ from $y$} if $x,y\notin A$ and $y\notin\comp Ax$.
\end{itemize}
\end{definition}

\formal{\lean{HPC.Plane.SeparatesPlane}, \lean{HPC.Plane.SeparatesPts} (\lean{connectedComponentIn}
for $\comp Ax$).}

If $A$ is closed, the components of $\R^2\setminus A$ are open (Lemma~\ref{lem:frontier-component}),
and an open connected subset of $\R^2$ is path-connected (the points joined to a given point by a
path inside the set form a clopen subset of it, since balls are convex). Separation is monotone in
the following sense: if $A\subseteq A'$, $x,y\notin A'$ and $A$ separates $x$ from $y$, then so does
$A'$ (since $\comp{A'}x\subseteq\comp Ax$).

\formal{\lean{IsOpen.isConnected_iff_isPathConnected} (Mathlib); \lean{HPC.Plane.separatesPts_mono}.}

\begin{lemma}\label{lem:compl-square}
For every $c\in\R^2$ and $R\ge0$ the complement of the closed ball $\bar B(c,R)$ (a closed square) is
connected and unbounded.
\end{lemma}

\begin{proof}
It is the union of the four open half-planes $\set{z_1>c_1+R}$, $\set{z_1<c_1-R}$,
$\set{z_2>c_2+R}$, $\set{z_2<c_2-R}$; each is convex, and each meets the next two (e.g.\ the first
and the third meet near $(c_1+R+1,c_2+R+1)$). So the union is connected; it is clearly unbounded.
\end{proof}

\formal{\lean{HPC.Plane.isConnected_compl_closedBall_pt},
\lean{HPC.Plane.isPreconnected_compl_closedBall}.}

\begin{lemma}\label{lem:unbounded-comp}
Let $K\subseteq\R^2$ be compact.
\begin{enumerate}[label=(\alph*)]
\item The complement of $K$ has at most one unbounded component; it contains the complement of
every closed ball $\bar B(0,R)\supseteq K$.
\item If $K$ separates the plane, some point $p\notin K$ has a bounded component $\comp Kp$.
\end{enumerate}
\end{lemma}

\begin{proof}
(a) Let $\bar B(0,R)\supseteq K$. The connected set $\R^2\setminus\bar B(0,R)$ misses $K$, so it lies
in a single component $C_\infty$. An unbounded component is not contained in $\bar B(0,R)$, so it
meets $\R^2\setminus\bar B(0,R)\subseteq C_\infty$ and equals $C_\infty$.
(b) $\R^2\setminus K$ is nonempty. If all its components were unbounded they would all equal
$C_\infty$ by (a), and $\R^2\setminus K$ would be connected.
\end{proof}

\formal{(a) \lean{HPC.Moore.mem_component_of_unbounded}; (b)
\lean{HPC.Plane.exists_bounded_component_of_separates}.}

\section{The two basic theorems}

We write the elements of $\R^2$ as complex numbers (Section~\ref{sec:intro-conv}); so $z-x$ is a
complex number and $\abs{z-x}$ its modulus.

\begin{theorem}[Same component $\Rightarrow$ logarithm]\label{thm:log-same-comp}
Let $K\subseteq\R^2$ be closed, $x\notin K$ and $y\in\comp Kx$. Then $z\mapsto\frac{z-x}{z-y}$ has
a logarithm on $K$.
\end{theorem}

\begin{proof}
$\comp Kx$ is open and connected, hence path-connected; let $\gamma\colon[0,1]\to\comp Kx$ be a path
from $y$ to $x$. The function $H(t,z)=\frac{z-\gamma(t)}{z-y}$ is continuous on $[0,1]\times K$ and
never zero ($\gamma(t)\notin K$ and $y\notin K$), and $H(0,z)=1$. By
Corollary~\ref{cor:log-exist}(c), $z\mapsto H(1,z)=\frac{z-x}{z-y}$ has a logarithm on $K$.
\end{proof}

\formal{\lean{HPC.Moore.exists_log_of_mem_component}.}

The next theorem is the converse direction, in a form with powers and an extra factor that is
needed for the invariance theorem.

\begin{theorem}[Bounded component $\Rightarrow$ no logarithm]\label{thm:nolog-bounded}
Let $K\subseteq\R^2$ be compact, $x\notin K$, and suppose $U=\comp Kx$ is bounded. Let $a\neq0$ be an
integer and $u$ a function continuous and nowhere zero on $\cl U$. Then
$z\mapsto(z-x)^a\,u(z)$ has no logarithm on $K$.
\end{theorem}

\begin{proof}
Suppose $g$ is one. By Tietze's extension theorem $g$ extends to a continuous $G\colon\R^2\to\C$.
The frontier of $U$ lies in $K$ (Lemma~\ref{lem:frontier-component}). Define
\[
\Psi(z)=\begin{cases} e^{G(z)}/u(z), & z\in U,\\ (z-x)^a, & z\notin U.\end{cases}
\]
The first formula is continuous on $\cl U$, the second on $\R^2\setminus U$ (which is closed and
does not contain $x\in U$), and they agree on $\cl U\setminus U=\operatorname{fr}U\subseteq K$, where
$e^{G}=e^{g}=(z-x)^au$. Hence $\Psi$ is continuous on $\R^2$; it vanishes nowhere. By
Corollary~\ref{cor:log-exist}(a) it has a logarithm $L$ on $\R^2$.

Choose $R_0$ with $U\subseteq\bar B(x,R_0)$, put $R=2\abs{R_0}+2$, and consider the euclidean circle
$c(t)=x+Re^{it}$, $t\in[0,2\pi]$. Its points have $\nrm{c(t)-x}\ge R/2>R_0$ (since $\abs w\le2\nrm
w$), so $c(t)\notin U$ and $\Psi(c(t))=(Re^{it})^a=e^{a(\ln R+it)}$. The functions $t\mapsto L(c(t))$
and $t\mapsto a(\ln R+it)$ are continuous on $[0,2\pi]$ with the same exponential, so by
Lemma~\ref{lem:log-unique} their difference is constant. Since $c(0)=c(2\pi)$, comparing $t=0$ and
$t=2\pi$ gives $2\pi i a=0$, a contradiction.
\end{proof}

\formal{\lean{HPC.Moore.no_log_pow_of_bounded_component}.}

\begin{corollary}\label{cor:nolog-bounded}
Let $K$ be compact, $y\notin K$, and suppose $\comp Kx$ is bounded and does not contain $y$. Then
$\frac{z-x}{z-y}$ has no logarithm on $K$. If $K$ separates the plane, there is $p\notin K$ such that
$z-p$ has no logarithm on $K$.
\end{corollary}

\begin{proof}
($x\notin K$, for otherwise the function vanishes at $x\in K$.) The closure of $U=\comp Kx$ lies in
$U\cup K$, which does not contain $y$; so $u(z)=1/(z-y)$ is continuous and nonzero on $\cl U$, and
Theorem~\ref{thm:nolog-bounded} with $a=1$ applies. For the second statement take $p$ from
Lemma~\ref{lem:unbounded-comp}(b), $a=1$ and $u=1$.
\end{proof}

\formal{\lean{HPC.Moore.no_log_of_bounded_component} (proved directly, by the same argument).}

\section{Janiszewski's theorem}

\begin{theorem}[Janiszewski]\label{thm:janiszewski}
Let $A,B\subseteq\R^2$ be compact with $A\cap B$ connected or empty. If neither $A$ nor $B$
separates $x$ from $y$, then $A\cup B$ does not separate $x$ from $y$.
\end{theorem}

\begin{proof}
Suppose $A\cup B$ separates $x$ from $y$; then $x,y\notin A\cup B$. As $A$ does not separate them,
$y\in\comp Ax$, and by Theorem~\ref{thm:log-same-comp} there is a logarithm $g_A$ of
$f(z)=\frac{z-x}{z-y}$ on $A$; likewise $g_B$ on $B$. On $A\cap B$ both are logarithms of $f$, so
if $A\cap B$ is connected, $g_A-g_B$ is a constant $c$ there (Lemma~\ref{lem:log-unique}), with
$e^c=1$; if $A\cap B=\emptyset$ put $c=0$. Define $g=g_A$ on $A$ and $g=g_B+c$ on $B\setminus A$.
The two formulas agree on $A\cap B$, and $A,B$ are closed, so $g$ is continuous on $A\cup B$; it is a
logarithm of $f$ there.

If $\comp{A\cup B}x$ is bounded, Corollary~\ref{cor:nolog-bounded} (with $K=A\cup B$; it does not
contain $y$) forbids $g$. If $\comp{A\cup B}y$ is bounded, the same corollary with $x$ and $y$
exchanged forbids the logarithm $-g$ of $\frac{z-y}{z-x}$. If both are unbounded they coincide by
Lemma~\ref{lem:unbounded-comp}(a), contradicting the assumption that $A\cup B$ separates $x$ from
$y$.
\end{proof}

\formal{\lean{HPC.Moore.janiszewski_compact}.}

A \emph{polygonal set} is a finite union of closed segments (degenerate segments, that is points,
are allowed); polygonal sets are compact.

\begin{corollary}[Janiszewski for polygonal sets]\label{cor:janiszewski-poly}
Theorem~\ref{thm:janiszewski} holds for polygonal sets $A$, $B$.
\end{corollary}

\formal{\lean{HPC.Plane.janiszewski_polygonal_of_compact}. The formal development also contains an
independent combinatorial proof, \lean{HPC.Plane.janiszewski_polygonal}, by mod-2 cycles of
straight-line graphs; see Remark~\ref{rem:jan-cycles}.}

\begin{remark}
Theorem~\ref{thm:janiszewski} fails for closed sets that are not compact: the two half-lines
$[0,\infty)\times\set0$ and $(-\infty,0]\times\set0$ meet in one point, neither separates $(0,1)$
from $(0,-1)$, and their union, the $x$-axis, does. (\lean{HPC.Plane.janiszewski_closed_false}.)
\end{remark}

\section{Borsuk's theorem}

The following theorem is the half of Borsuk's characterization of plane separation used later.

\begin{theorem}[Borsuk]\label{thm:borsuk}
Let $K\subseteq\R^2$ be compact with $\R^2\setminus K$ connected. Every function $f$ continuous and
nowhere zero on $K$ has a logarithm on $K$.
\end{theorem}

The idea: approximate $f$ by a smooth map $\Phi$ whose zeros are finitely many nondegenerate points
outside $K$; divide out the zeros one at a time; each factor $z-z_0$ (or its conjugate) has a
logarithm on $K$ because $z_0$ can be joined to infinity in $\R^2\setminus K$.

\begin{definition}\label{def:good-zero}
Write $\sigma_+(w)=w$ and $\sigma_-(w)=\cl w$. A \emph{good zero} of a function $\Phi\colon\R^2\to\C$
is a point $z_0$ for which there are $r>0$, a sign $s\in\set{+,-}$, $\gamma\in\C$ and $0\le\rho<\abs\gamma$
with
\[
\abs{\Phi(z)-\gamma\,\sigma_s(z-z_0)}\le\rho\abs{z-z_0}\qquad\text{whenever } d(z,z_0)<r.
\]
\end{definition}

\formal{\lean{HPC.Moore.IsGoodZero}, \lean{HPC.Moore.sgnConj}.}

Taking $z=z_0$ shows $\Phi(z_0)=0$.

\begin{lemma}\label{lem:good-zero}
\leavevmode
\begin{enumerate}[label=(\alph*)]
\item A good zero is isolated: $\Phi$ has no other zero in $B(z_0,r)$.
\item If $G\colon\R^2\to\R^2$ is differentiable at $z_0$ with $G(z_0)=c$ and $\det DG(z_0)\neq0$,
then $z_0$ is a good zero of $\Phi(z)=G(z)-c$ (as a complex number).
\item If $z_0$ is a good zero of $\Phi$, $m$ is continuous at $z_0$ with $m(z_0)\neq0$ and
$\Psi=\Phi m$ near $z_0$, then $z_0$ is a good zero of $\Psi$.
\item A continuous function whose zeros in a compact set $B$ are all good has finitely many zeros
in $B$.
\end{enumerate}
\end{lemma}

\begin{proof}
(a) If $\Phi(z)=0$ with $0<d(z,z_0)<r$, then $\abs\gamma\abs{z-z_0}\le\rho\abs{z-z_0}$, so
$\abs\gamma\le\rho$.

(b) Write $A=DG(z_0)$. Every real-linear map of $\C$ has the form $Aw=\alpha w+\beta\cl w$, and
$\det A=\abs\alpha^2-\abs\beta^2$ (a direct computation with the matrix of $A$). So
$\abs\alpha\neq\abs\beta$. By differentiability, for every $\eta>0$ there is $r>0$ with
$\abs{G(z)-c-A(z-z_0)}\le\eta\abs{z-z_0}$ for $d(z,z_0)<r$. If $\abs\alpha>\abs\beta$, take $s=+$,
$\gamma=\alpha$ and $\eta=(\abs\alpha-\abs\beta)/4$: then
$\abs{\Phi(z)-\alpha(z-z_0)}\le(\abs\beta+\eta)\abs{z-z_0}$ and $\rho=\abs\beta+\eta<\abs\alpha$. If
$\abs\beta>\abs\alpha$, take $s=-$ and $\gamma=\beta$ symmetrically.

(c) Let $m_0=m(z_0)$ and choose $\delta>0$ small (explicitly
$\delta=(\abs\gamma-\rho)\abs{m_0}/2(\rho+\abs\gamma+1)$) and $r'>0$ so that for $d(z,z_0)<r'$ we
have $d(z,z_0)<r$, $\abs{m(z)-m_0}<\delta$ and $\Psi(z)=\Phi(z)m(z)$. With $w=z-z_0$,
\[
\Psi(z)-\gamma m_0\sigma_s(w)=(\Phi(z)-\gamma\sigma_s(w))m(z)+\gamma\sigma_s(w)(m(z)-m_0),
\]
so $\abs{\Psi(z)-\gamma m_0\sigma_s(w)}\le\big(\rho(\abs{m_0}+\delta)+\abs\gamma\delta\big)\abs w$, and
the choice of $\delta$ makes $\rho(\abs{m_0}+\delta)+\abs\gamma\delta<\abs{\gamma m_0}$.

(d) Otherwise the zeros in $B$ form an infinite closed subset of the compact set $B$, which has an
accumulation point; it is a zero, and by (a) not an accumulation point of zeros.
\end{proof}

\formal{(a) \lean{HPC.Moore.IsGoodZero.isolated}; (b) \lean{HPC.Moore.isGoodZero_of_hasFDerivAt};
(c) \lean{HPC.Moore.IsGoodZero.mul}; (d) \lean{HPC.Moore.finite_zeros_of_goodZeros}.}

\begin{lemma}[Removing a good zero]\label{lem:remove-zero}
Let $\Phi$ be continuous, $z_0$ a good zero of $\Phi$ and $r>0$. There are a sign $s$ and a
continuous $\Phi'$ with
\begin{enumerate}[label=(\roman*)]
\item $\Phi(z)=\Phi'(z)\,\sigma_s(z-z_0)$ whenever $d(z,z_0)\ge r$;
\item $\Phi'(z)=0$ if and only if $z\neq z_0$ and $\Phi(z)=0$;
\item every good zero $z\neq z_0$ of $\Phi$ is a good zero of $\Phi'$.
\end{enumerate}
\end{lemma}

\begin{proof}
Let $r_0,s,\gamma,\rho$ be as in Definition~\ref{def:good-zero}; $\gamma\neq0$. Put
$r_1=\min(r,r_0/2)$, $\psi(z)=\min\big(1,d(z,z_0)/r_1\big)$ and, with $w=z-z_0$,
\[
F(z)=\gamma+\psi(z)\Big(\frac{\Phi(z)}{\sigma_s(w)}-\gamma\Big)\quad(z\neq z_0),\qquad
\Phi'(z_0)=\gamma,\quad\Phi'(z)=F(z)\ (z\neq z_0).
\]
For $0<d(z,z_0)<r_0$ the good-zero estimate divided by $\abs w$ gives
$\abs{\Phi(z)/\sigma_s(w)-\gamma}\le\rho$, hence $\abs{F(z)-\gamma}\le\psi(z)\rho\le\rho\, d(z,z_0)/r_1$.
This shows that $\Phi'$ is continuous at $z_0$, and that $\Phi'$ has no zero with
$d(z,z_0)<r_1$ (there $\abs{F-\gamma}\le\rho<\abs\gamma$); and $\Phi$ has none either for
$0<d(z,z_0)<r_0$. For $d(z,z_0)\ge r_1$ we have $\psi=1$ and $F=\Phi/\sigma_s(w)$, which gives (i) and,
together with the previous sentence, (ii). Away from $z_0$, $\Phi'$ is continuous as a composition.
For (iii): a zero $z\neq z_0$ of $\Phi$ has $d(z,z_0)\ge r_0>r_1$, and near it
$\Phi'=\Phi\cdot(1/\sigma_s(w))$; apply Lemma~\ref{lem:good-zero}(c).
\end{proof}

\formal{\lean{HPC.Moore.IsGoodZero.remove}.}

\begin{lemma}\label{lem:log-factor}
Let $K$ be compact with $\R^2\setminus K$ connected, $z_0\notin K$ and $s$ a sign. Then
$\sigma_s(z-z_0)$ has a logarithm on $K$.
\end{lemma}

\begin{proof}
Choose $R$ with $K\subseteq\bar B(0,R)$ and the far point $y=(\abs R+1,0)$; then $y\notin K$ and
$\Re(y-z)>0$ for $z\in K$. Since $\R^2\setminus K$ is connected, $y\in\comp K{z_0}$, and
Theorem~\ref{thm:log-same-comp} gives a logarithm $g_1$ of $\frac{z-z_0}{z-y}$ on $K$. Now
$\Log(y-z)+i\pi$ is a logarithm of $z-y$ on $K$, so $g=g_1+\Log(y-z)+i\pi$ is a logarithm of $z-z_0$.
For $s=-$ use $\cl g$.
\end{proof}

\formal{\lean{HPC.Moore.exists_log_sgnConj}.}

\begin{lemma}[Induction on the zeros]\label{lem:zeros-induction}
Let $K$ be compact with $\R^2\setminus K$ connected, and $B\supseteq K$ convex. Let $\Phi\colon\R^2\to\C$
be continuous, nowhere zero on $K$, with finitely many zeros in $B$, all of them good. Then $\Phi$
has a logarithm on $K$.
\end{lemma}

\begin{proof}
Induction on the number $n$ of zeros in $B$ (we may assume $K\neq\emptyset$). If $n=0$, $\Phi$ has a
logarithm on $B$ by Corollary~\ref{cor:log-exist}(b), and we restrict it to $K$. If $n>0$, pick a
zero $z_0\in B$; $z_0\notin K$, so $r=\dist(z_0,K)>0$. Lemma~\ref{lem:remove-zero} gives $s$ and
$\Phi'$; by (ii) and (iii), $\Phi'$ has the $n-1$ zeros of $\Phi$ in $B$ other than $z_0$, all good,
and none on $K$. By induction $\Phi'$ has a logarithm on $K$, and so has $\sigma_s(z-z_0)$
(Lemma~\ref{lem:log-factor}). Every $z\in K$ has $d(z,z_0)\ge r$, so $\Phi=\Phi'\sigma_s(z-z_0)$ on $K$
by (i), and the sum of the two logarithms is one of $\Phi$.
\end{proof}

\formal{\lean{HPC.Moore.exists_log_of_goodZeros}.}

Two standard theorems of analysis enter the proof of Theorem~\ref{thm:borsuk}; both are quoted.
\begin{itemize}
\item \emph{Stone--Weierstrass.} The real $C^1$ functions on $\R^2$ form an algebra of continuous
functions containing the constants and separating points (the coordinate functions do); hence every
continuous real function on a compact set $K$ is a uniform limit on $K$ of $C^1$ functions.
\item \emph{Sard's theorem in equal dimensions.} If $G\colon\R^2\to\R^2$ is differentiable, the image
of the set $\set{z:\det DG(z)=0}$ has Lebesgue measure zero. In particular every ball contains a
\emph{regular value} $c$: a point such that $\det DG(z)\neq0$ whenever $G(z)=c$.
\end{itemize}

\formal{Mathlib's
\lean{ContinuousMap.exists_mem_subalgebra_near_continuous_of_isCompact_of_separatesPoints} and
\lean{MeasureTheory.addHaar_image_eq_zero_of_det_fderivWithin_eq_zero}; in the development
\lean{HPC.Moore.exists_contDiff_near} and \lean{HPC.Moore.exists_small_regular_value}.}

\begin{proof}[Proof of Theorem~\ref{thm:borsuk}]
We may assume $K\neq\emptyset$. Let $m=\min_K\abs f>0$. Extend $f$ continuously to $\R^2$ (Tietze)
and approximate its real and imaginary parts on $K$ by $C^1$ functions: there is a $C^1$ map
$G\colon\R^2\to\R^2$ with $\abs{G(z)-f(z)}<m/4$ on $K$. Choose a regular value $c$ of $G$ with
$\nrm c<m/8$, so $\abs c<m/4$, and put $\Phi(z)=G(z)-c$. On $K$,
$\abs{\Phi-f}<m/2<\abs f$; in particular $\Phi$ has no zero on $K$. Every zero $z$ of $\Phi$ has
$G(z)=c$, hence $\det DG(z)\neq0$, and is a good zero by Lemma~\ref{lem:good-zero}(b). Let $B$ be a
closed ball containing $K$; it is convex and compact, so $\Phi$ has finitely many zeros in $B$
(Lemma~\ref{lem:good-zero}(d)). By Lemma~\ref{lem:zeros-induction} $\Phi$ has a logarithm on $K$, and
by Lemma~\ref{lem:log-near} so has $f$.
\end{proof}

\formal{\lean{HPC.Moore.borsuk_exists_log}.}

\section{Plane separation is a topological invariant}

\begin{theorem}\label{thm:sep-invariant}
Let $A,B\subseteq\R^2$ be compact and homeomorphic. If $A$ separates the plane, so does $B$.
\end{theorem}

\begin{proof}
Let $\varphi\colon A\to B$ be a homeomorphism and, by Lemma~\ref{lem:unbounded-comp}(b), $x\notin A$
with $\comp Ax$ bounded. The function $F(z)=\varphi^{-1}(z)-x$ is continuous and nowhere zero on $B$.
If $B$ did not separate the plane, Theorem~\ref{thm:borsuk} would give a logarithm $g$ of $F$ on
$B$, and then $g\circ\varphi$ would be a logarithm of $z-x$ on $A$, contradicting
Theorem~\ref{thm:nolog-bounded} (with $a=1$, $u=1$).
\end{proof}

\formal{\lean{HPC.Jones.separatesPlane_of_homeomorph}.}

\begin{remark}
A weaker form of this invariance, proved without Borsuk's theorem, appears in
Section~\ref{sec:weak-invariance}: if $K$ separates the plane and the continuum $K'$ is homeomorphic
to $K$, some subcontinuum of $K'$ separates the plane. The formal development uses it in Hagopian's
argument (Volume~6); Theorem~\ref{thm:sep-invariant} is used for Jones's decomposition (Volume~7).
\end{remark}

\section*{Exercises}

\begin{exercise}
Let $K$ be compact and $x,y\notin K$ in the same component of $\R^2\setminus K$. Show that the
logarithm of $\frac{z-x}{z-y}$ given by Theorem~\ref{thm:log-same-comp} can be chosen to tend to
$0$ as $x\to y$. (Use the homotopy along a short segment.)
\end{exercise}

\begin{exercise}[Mullikin]
Let $M_1,M_2$ be continua with $M_1\cap M_2=K_1\cup K_2$ for disjoint nonempty closed sets
$K_1,K_2$. Show that $M_1\cup M_2$ separates the plane. (Let $\varphi\colon M_2\to[0,1]$ be continuous,
$0$ on $K_1$ and $1$ on $K_2$. Show that the function equal to $1$ on $M_1$ and to $e^{2\pi i\varphi}$
on $M_2$ has no logarithm on $M_1\cup M_2$, using Lemma~\ref{lem:int-valued} on $M_1$ and on $M_2$;
then apply Theorem~\ref{thm:borsuk}. This is \lean{HPC.Moore.not_isConnected_compl_union}. If
moreover $K_1$, $K_2$ are connected and neither $M_1$ nor $M_2$ separates the plane, the complement
has exactly two components: \lean{HPC.Moore.mullikin_theorem5}.)
\end{exercise}

\begin{exercise}
Show that the hypothesis ``$\R^2\setminus K$ connected'' cannot be dropped from
Theorem~\ref{thm:borsuk}: take for $K$ the unit circle and $f(z)=z$. Which step of the proof
fails?
\end{exercise}
